\section*{Capítulo \ref{ch:b2:magmatic-differentiation} --- La diferenciación magmática}

\begin{solution}{exo:b2:magmatic-differentiation:1}
(a) El liquidus de An$_{50}$ está a \qty{1450}{\celsius}. El primer cristal está sobre el
solidus a esa temperatura, entre An$_{66.7}$ (1372) y An$_{83.3}$ (1465):
$66.7 + 16.6\times(1450 - 1372)/(1465 - 1372) = 66.7 + 13.9 = \qty{80.6}{\percent}$
de anortita. (b) A \qty{1394}{\celsius} el líquido está sobre el liquidus,
An$_{33.3}$; los cristales sobre el solidus, $66.7 + 16.6\times22/93 =
\qty{70.6}{\percent}$. Por la regla de la palanca, la fracción cristalizada es $(50 -
33.3)/(70.6 - 33.3) = 16.7/37.3 = 0.448$. (c) El último líquido desaparece cuando los
cristales alcanzan An$_{50}$, a la temperatura del solidus, \qty{1287}{\celsius}; el
liquidus está allí entre An$_{11.1}$ (1265) y An$_{20}$ (1334): $11.1 +
8.9\times22/69 = \qty{13.9}{\percent}$ de anortita.
\end{solution}

\begin{solution}{exo:b2:magmatic-differentiation:2}
$\bar D_{\ce{Sr}} = 0.15\times0.000063 + 0.35\times0.157 + 0.50\times2.7 = 0.00001 +
0.0550 + 1.35 = 1.40$; $\bar D_{\ce{Zr}} = 0.0001 + 0.0683 + 0.0005 = 0.069$;
$\bar D_{\ce{Rb}} = 0.00001 + 0.0012 + 0.0125 = 0.014$. Solo el estroncio, retenido por la
plagioclasa, es compatible ($\bar D > 1$); el circonio y el rubidio son incompatibles,
el rubidio más aún.
\end{solution}

\begin{solution}{exo:b2:magmatic-differentiation:3}
(a) $C_L/C_0 = 0.5^{D-1}$: Sr $0.5^{0.40} = 0.758$ (disminuye); Zr $0.5^{-0.931} =
1.907$; Rb $0.5^{-0.986} = 1.981$, cerca del límite 2. (b) Zr/Rb cambia en
$0.5^{0.069 - 0.014} = 0.5^{0.055} = 0.963$, menos de un \qty{4}{\percent} mientras cada
elemento casi se duplicaba: como dice el corolario, la razón de dos \omterm{def:b2:elements-in-the-earth:partition}{elementos incompatibles}
conserva el valor del \omterm{def:b2:magma-emplacement:magma}{magma} parental, y por tanto informa sobre la fuente, no
sobre la cristalización.
\end{solution}

\begin{solution}{exo:b2:magmatic-differentiation:4}
(a) \ce{SiO2} $+5.03$, \ce{MgO} $-1.95$, \ce{K2O} $+0.82$ (por ciento); la potasa
se multiplica por $1.67/0.85 = 1.96$. (b) Para un elemento perfectamente incompatible
$C_L/C_0 = 1/F$, luego $F = 0.85/1.67 = 0.509$: tendría que haber cristalizado la mitad del
\omterm{def:b2:magma-emplacement:magma}{magma}. (c) Retirar \qty{12.5}{g} por cada \qty{100}{g} deja $F \approx 0.88$ y
solo elevaría la potasa en un factor de $1/0.88 = 1.14$: la cristalización sola no puede
explicar el aumento. La potasa vino de fuera, de los xenolitos graníticos, ricos en
potasa (\qty{3.6}{\percent}), que el \omterm{def:b2:magma-emplacement:magma}{magma} asimiló.
\end{solution}

\begin{solution}{exo:b2:magmatic-differentiation:5}
(a) $F = 0.8$, $D = 5$. Fraccionada: $0.8^{4} = 0.410$; en equilibrio:
$1/(5 + 0.8\times(1 - 5)) = 1/1.8 = 0.556$. (b) $F = 0.1$, $D = 0.01$. Fraccionada:
$0.1^{-0.99} = 9.77$; en equilibrio: $1/(0.01 + 0.1\times0.99) = 1/0.109 = 9.17$.
(c) Cuando $F \to 0$, la ley de Rayleigh tiende a $0$ para $D > 1$ y a infinito para
$D < 1$, mientras que la ley de equilibrio tiende a $1/D$ (el último líquido en equilibrio
con cristales de la composición global). Solo la \omterm{def:b2:magmatic-differentiation:crystallisation}{cristalización fraccionada} puede despojar a un
líquido de un \omterm{def:b2:elements-in-the-earth:partition}{elemento compatible} por completo.
\end{solution}

\begin{solution}{exo:b2:magmatic-differentiation:6}
(a) Una unidad fórmula contiene 1.6 \ce{MgO}, 0.4 \ce{FeO} y un \ce{SiO2}: $1.6\times
40.30 = 64.48$, $0.4\times71.84 = 28.74$, 60.08, en total \qty{153.30}{g}. Porcentajes
en masa: \ce{MgO} \qty{42.06}{\percent}, \ce{FeO} \qty{18.75}{\percent},
\ce{SiO2} \qty{39.19}{\percent}. (b) Regla de la palanca sobre \ce{MgO}: $x = (10.0 - 8.0)/(42.06
- 8.0) = 2.0/34.06 = 0.0587$, cerca de un \qty{6}{\percent} de olivino. (c) Sílice del
líquido: $(49.0 - 0.0587\times39.19)/(1 - 0.0587) = 46.70/0.9413 =
\qty{49.61}{\percent}$. Retirar olivino, más pobre en sílice que el basalto, eleva la
sílice del líquido, pero solo despacio.
\end{solution}

\begin{solution}{exo:b2:magmatic-differentiation:7}
(a) $F(0) = 0.768$; $F(1000) = 0.1091 - 0.5906 + 0.7678 = 0.286$; $F(2000) = 0.4364
- 1.1813 + 0.7678 = 0.023$. (b) El discriminante es $(5.9064\times10^{-4})^2 -
4\times1.091\times10^{-7}\times0.7678 = 3.4885\times10^{-7} - 3.3507\times10^{-7} =
1.379\times10^{-8}$, de raíz cuadrada $1.174\times10^{-4}$; la raíz menor es
$H = (5.9064 - 1.174)\times10^{-4}/(2.182\times10^{-7}) = \qty{2169}{m}$, el
Horizonte Sándwich. (c) $1 - 0.768 = 0.232$: casi una cuarta parte del \omterm{def:b2:magma-emplacement:magma}{magma} había
cristalizado antes de que se formara la base de la Zona Inferior aflorante, en la Zona Oculta
bajo ella y en las Series de Borde a lo largo de las paredes y el techo.
\end{solution}

\begin{solution}{exo:b2:magmatic-differentiation:8}
El \omterm{def:b2:magmatic-differentiation:cumulate}{cumulado} contiene $D\,C_L = D\,C_0F^{D-1}$, de modo que entre dos niveles
$\ln(\mathrm{Ni}_2/\mathrm{Ni}_1) = (D - 1)\ln(F_2/F_1)$. (a) $D - 1 = \ln(54/177)/
\ln(0.29/0.70) = (-1.187)/(-0.881) = 1.35$, luego $D \approx 2.35$ (la regresión sobre
los 66 \omterm{def:b2:magmatic-differentiation:cumulate}{cumulados} con $F \ge 0.23$ da 2.38). (b) $D - 1 = \ln(6.2/54)/\ln(0.17/0.29) = (-2.164)/(-0.534)
= 4.05$, luego $D \approx 5.05$. (c) El $D$ global cambia cuando cambia el conjunto
que cristaliza (un mineral nuevo que capta níquel, o uno que deja de cristalizar), y
los propios coeficientes mineral--fundido aumentan a medida que el líquido se enfría y cambia de
composición, como se ha medido para el olivino; el líquido atrapado en los \omterm{def:b2:magmatic-differentiation:cumulate}{cumulados} difumina la
estimación.
\end{solution}

\begin{solution}{exo:b2:magmatic-differentiation:9}
(a) $100 + 25.4 - 2.9 - 9.6 = \qty{112.9}{g}$. (b) Gramos de \ce{SiO2}: $55.53 +
25.4\times0.7248 - 2.9\times0.3919 - 9.6\times0.5054 = 55.53 + 18.41 - 1.14 - 4.85 =
\qty{67.95}{g}$, luego $67.95/112.9 = \qty{60.19}{\percent}$ frente a 60.16. \ce{MgO}:
$4.66 + 0.20 - 1.22 = \qty{3.64}{g}$, \qty{3.22}{\percent} frente a 3.19. \ce{K2O}:
$0.88 + 0.92 = \qty{1.80}{g}$, \qty{1.60}{\percent} frente a 1.75. (c) La sílice y la magnesia, dos de los tres
controles, se ajustan con un margen de 0.03; la potasa, que no se usó como control, queda 0.15 por debajo, casi una
décima parte de su valor: el material asimilado tenía más potasa que el xenolito
medio, o se incorporó preferentemente un fundido de él rico en potasa.
\end{solution}

\begin{solution}{exo:b2:magmatic-differentiation:10}
(a) $z = (0.2 + 1.5 - 1)/(0.2 - 1) = 0.7/(-0.8) = -0.875$; $F^{-z} = 0.7^{0.875} =
0.7319$; $C^*/C_0 = 0.2\times0.5/(0.2 - 1 + 1.5) = 0.1/0.7 = 0.1429$, luego $C/C_0 =
0.7319 + 0.1429\times(1 - 0.7319) = 0.7319 + 0.0383 = 0.770$. (b) Rayleigh:
$0.7^{0.5} = 0.837$. El \omterm{def:b2:magma-emplacement:magma}{magma} con AFC tiene menos estroncio por dos razones: para perder
un \qty{30}{\percent} de su masa ganando $r$ veces lo que cristaliza debe
cristalizar $(1 - 0.7)/(1 - 0.2) = 0.375$ en lugar de 0.30, y el estroncio que gana
procede de una roca la mitad de rica. (c) $(0.7319\times0.7035 + 0.0383\times0.7200)/0.7702
= (0.5149 + 0.0276)/0.7702 = 0.70432$: la razón ha subido 0.0008, mientras que la
cristalización sola la habría dejado en 0.7035.
\end{solution}

\begin{solution}{exo:b2:magmatic-differentiation:11}
(a) $1.1\times(1100 - 200) + 250 = 990 + 250 = \qty{1240}{J}$ por gramo. (b)
$25.4\times1240 = \qty{31496}{J}$, cerca de los \qty{31350}{J} de Wilcox. (c)
$31496/370 = \qty{85.1}{g}$ de plagioclasa. (d) El balance de masas solo retiró $2.9 +
9.6 = \qty{12.5}{g}$ de cristales, cuyo calor latente (unos \qty{4600}{J}) cubre menos
de una sexta parte de lo necesario; $r = 25.4/12.5 = 2.03$, el doble de roca asimilada que de
cristales formados, mientras que el balance térmico solo permite un $r$ muy inferior a 1. El \omterm{def:b2:magma-emplacement:magma}{magma} emitido
no pudo aportar el calor por sí mismo: una masa mayor de \omterm{def:b2:magma-emplacement:magma}{magma} debió de cristalizar
y enfriarse en profundidad, fuera de la vista.
\end{solution}

\begin{solution}{exo:b2:magmatic-differentiation:12}
(a) Etapa 1, $F$ de 1 a 0.8 con $D = 5$: factor $0.8^{4} = 0.4096$. Etapa 2,
de 0.8 a 0.5 con $D = 1.5$: la ley se aplica al líquido restante, así que el
factor es $(0.5/0.8)^{0.5} = 0.7906$. Níquel: $250\times0.4096\times0.7906 =
\qty{81.0}{ppm}$. (b) $250\times0.5^{4} = \qty{15.6}{ppm}$, cinco veces demasiado bajo. (c)
La ley de Rayleigh supone un $D$ constante; cuando cambia el conjunto que cristaliza,
$D$ salta, y la ley debe reiniciarse desde la composición del líquido alcanzada en el
cambio, con $F$ medido respecto de ese líquido.
\end{solution}

\begin{solution}{pb:b2:magmatic-differentiation:1}
\noindent\textbf{Parte I.}

\textbf{1.} $\bar D_{\ce{Ni}} = 0.20\times5.0 + 0.50\times0 + 0.30\times2.0 = 1.60$.

\textbf{2.} $\bar D_{\ce{Zr}} = 0.20\times0.0007 + 0.50\times0.001 + 0.30\times0.195 =
0.00014 + 0.0005 + 0.0585 = 0.059$.

\textbf{3.} Zona superior: $\bar D_{\ce{Ni}} = 0.45\times2.0 = 0.90$; $\bar D_{\ce{Zr}} =
0.55\times0.001 + 0.45\times0.195 = 0.088$.

\textbf{4.} El níquel es compatible en la zona inferior (1.60) y se vuelve ligeramente
incompatible (0.90) cuando desaparece el olivino; el circonio es incompatible en ambas.

\textbf{5.} El clinopiroxeno: aporta $0.0585$ de los $0.059$ de la zona inferior, casi
todo, porque es el único mineral con un coeficiente lejos de cero.

\textbf{6.} $(C_{\ce{Ni}}/C_{\ce{Zr}})$ cambia en $F^{D_{\ce{Ni}} - D_{\ce{Zr}}}$;
por cada reducción a la mitad, $0.5^{1.60 - 0.059} = 0.5^{1.541} = 0.344$.

\noindent\textbf{Parte II.}

\textbf{7.} $1.60\times150 = \qty{240}{ppm}$.

\textbf{8.} El líquido tenía $190/1.60 = \qty{118.75}{ppm}$; $F = (118.75/150)^{1/0.6} =
0.7917^{1.667} = 0.68$.

\textbf{9.} $80\times0.68^{-0.941} = 80\times1.438 = \qty{115}{ppm}$.

\textbf{10.} Fraccionada: $150\times0.68^{0.6} = 150\times0.793 = \qty{119.0}{ppm}$;
en equilibrio: $150/(1.6 + 0.68\times(1 - 1.6)) = 150\times0.839 = \qty{125.8}{ppm}$.
La diferencia es solo de un \qty{6}{\percent} después de cristalizar un tercio del
\omterm{def:b2:magma-emplacement:magma}{magma}: para un elemento moderadamente compatible las dos leyes se mantienen próximas al principio.

\textbf{11.} $150\times0.55^{0.6} = \qty{104.8}{ppm}$.

\textbf{12.} Justo por debajo: $1.60\times104.8 = \qty{167.7}{ppm}$; justo por encima:
$0.90\times104.8 = \qty{94.3}{ppm}$. El perfil de níquel de los \omterm{def:b2:magmatic-differentiation:cumulate}{cumulados} cae
bruscamente un \qty{44}{\percent} en el nivel donde el olivino desaparece de las
capas: un escalón, no una curva suave.

\textbf{13.} Líquido: $104.8\times(0.40/0.55)^{0.90 - 1} = 104.8\times1.032 =
\qty{108.2}{ppm}$; \omterm{def:b2:magmatic-differentiation:cumulate}{cumulado} $0.90\times108.2 = \qty{97.4}{ppm}$. Con $\bar D < 1$
el níquel sube ahora despacio en el líquido y en los \omterm{def:b2:magmatic-differentiation:cumulate}{cumulados}: el perfil, tras su
escalón, asciende suavemente.

\noindent\textbf{Parte III.}

\textbf{14.} $x_{\mathrm{ol}}(m_{\mathrm{ol},j} - d_j) + x_{\mathrm{pl}}
(m_{\mathrm{pl},j} - d_j) + x_{\mathrm{cpx}}(m_{\mathrm{cpx},j} - d_j) = 100\,(p_j -
d_j)$.

\textbf{15.} \ce{MgO}: $37.16\,x_{\mathrm{ol}} - 4.9\,x_{\mathrm{pl}} +
11.1\,x_{\mathrm{cpx}} = 410$; \ce{CaO}: $-7.3\,x_{\mathrm{ol}} +
7.06\,x_{\mathrm{pl}} + 13.2\,x_{\mathrm{cpx}} = 420$; \ce{Al2O3}:
$-14.1\,x_{\mathrm{ol}} + 17.6\,x_{\mathrm{pl}} - 11.1\,x_{\mathrm{cpx}} = 190$.

\textbf{16.} Resolviendo (por eliminación o por la regla de Cramer): $x_{\mathrm{ol}} = 9.03$,
$x_{\mathrm{pl}} = 30.84$, $x_{\mathrm{cpx}} = 20.32$ gramos por cada \qty{100}{g}.

\textbf{17.} Cristales $9.03 + 30.84 + 20.32 = \qty{60.19}{g}$; líquido
\qty{39.81}{g}: $F = 0.398$.

\textbf{18.} $(4950 - 9.03\times39.19 - 30.84\times50.54 - 20.32\times51.5)/39.81 =
(4950 - 353.9 - 1558.7 - 1046.5)/39.81 = 1990.9/39.81 = \qty{50.0}{\percent}$, como en
el \omterm{def:b2:magma-emplacement:bodies}{dique}: el óxido no usado como control se reproduce.

\textbf{19.} Los cristales retirados son $9.03/60.19 = \qty{15}{\percent}$ de olivino,
un \qty{51}{\percent} de plagioclasa y un \qty{34}{\percent} de clinopiroxeno, entre las
modas de las dos zonas. Corresponden al conjunto inferior cristalizando de $F = 1$ a
0.55 ($0.45\times20 = 9.0$, $0.45\times50 = 22.5$, $0.45\times30 = 13.5$) seguido del
superior de 0.55 a 0.40 ($0.15\times55 = 8.25$, $0.15\times45 = 6.75$): 9.0,
30.75 y 20.25 gramos.

\textbf{20.} $0.30/0.75 = 0.40$.

\noindent\textbf{Parte IV.}

\textbf{21.} $0.55^{0.059 - 1} = 0.55^{-0.941} = 1.755$; $80\times1.755 =
\qty{140.4}{ppm}$.

\textbf{22.} $188/140.4 = 1.339 = (F/0.55)^{0.088 - 1}$, luego $F = 0.55\times1.339^{-1/
0.912} = 0.55\times0.726 = 0.399$.

\textbf{23.} $(188/80)^{1/(0.07 - 1)} = 2.35^{-1.075} = 0.399$. Para los \omterm{def:b2:elements-in-the-earth:partition}{elementos incompatibles}
$C_L/C_0 \approx 1/F$ sea cual sea el pequeño $D$: la respuesta apenas depende de
los detalles de $D$, lo que hace de los \omterm{def:b2:elements-in-the-earth:partition}{elementos incompatibles} los mejores indicadores de $F$.

\textbf{24.} La plagioclasa, que no contiene hierro, forma la mitad de los cristales, y el
piroxeno y el olivino juntos contienen menos hierro que magnesia: el líquido
se enriquece en hierro y se empobrece en magnesia sin ganar sílice. Este enriquecimiento
en hierro con la sílice casi constante es la huella de la serie toleítica del
capítulo siguiente.

\textbf{25.} Los cristales cúmulo puros contendrían $0.088\times188 = \qty{16.5}{ppm}$. Con una
fracción másica $t$ de líquido atrapado, $16.5(1 - t) + 188\,t = 47$, luego $t = (47 -
16.5)/(188 - 16.5) = 30.5/171.5 = 0.18$: cerca de un \qty{18}{\percent} de líquido atrapado.

\textbf{26.} Las tres estimaciones, elementos mayores (0.398), potasa (0.40) y
circonio (0.399), coinciden: cuando cristalizó la capa más alta, seguía líquido un
\textbf{\qty[text-series-to-math=true]{40}{\percent}} del \omterm{def:b2:magma-emplacement:magma}{magma}
original.
\end{solution}
